\subsection{Motivation: stochastic decision problems}

So far, most of the control problems considered in these notes have been written in the form of
a deterministic state-space model
\[
x_{k+1} = f(x_k,u_k),
\]
possibly with constraints and disturbances. This representation is extremely useful when the
state is continuous, the dynamics are known with sufficient accuracy, and the effect of the input
can be described by deterministic equations.\\

However, there are many systems for which this description is not the most natural one. A
typical example is the control of traffic lights at an intersection. The control input is discrete:
at each time we decide which light is red and which light is green. The state may also be
naturally discrete: for instance, the number of vehicles in each queue. Moreover, the evolution
is stochastic, because new cars arrive randomly. In such a setting, writing a smooth deterministic
state-space model is not very convenient.\\
Consider, for simplicity, a single queue with queue length $q_k$. At each time step, a new
vehicle arrives with probability $p$. If the light is red, no cars are served. If the light is green,
up to three cars can leave the queue. A simple model is therefore
\[
a_k =
\begin{cases}
1 & \text{with probability } p\\
0 & \text{with probability } 1-p
\end{cases}
\]
and
\[
q_{k+1}
=
\begin{cases}
q_k + a_k & \text{if the light is red}\\
\max(0,q_k-3) + a_k & \text{if the light is green}.
\end{cases}
\]
If the queue length is bounded by a maximum capacity $q_{\max}$, one can additionally saturate
the update at $q_{\max}$.\\
This is still a Markovian model: once the current queue length \(q_k\) and the current action
are known, the probability distribution of \(q_{k+1}\) is fully determined. The future does not
depend on the entire past history of the queue, but only on the present state and the action
chosen now. The difference with the deterministic models used before is that the next state is
not a single value, but a probability distribution over possible next states.\\
This motivates the introduction of Markov Decision Processes.

\subsection{Finite Markov Decision Processes}

A finite Markov Decision Process, or MDP, is a controlled stochastic dynamical system described
by the following objects:
\[
\cX = {1,\dots,N},
\qquad
\cU = {1,\dots,M},
\]
where $(\cX)$ is the finite set of states and $(\cU)$ is the finite set of actions.\\
The dynamics are described by transition probabilities
\[
P : \cX \times \cU \times \cX \to [0,1].
\]
We write
\[
P^u_{x,x'}
=
\Prob[x_{k+1}=x' \mid x_k=x,\ u_k=u].
\]
Thus \(P^u_{x,x'}\) is the probability of moving from state \(x\) to state \(x'\) when action
\(u\) is applied.\\
For every fixed state-action pair \((x,u)\), the transition probabilities satisfy
\[
P^u_{x,x'} \ge 0,
\qquad
\sum_{x'\in\cX} P^u_{x,x'} = 1.
\]
The Markov property means precisely that
\[
\Prob[x_{k+1}=x' \mid x_k,u_k,x_{k-1},u_{k-1},\dots,x_0]
=
\Prob[x_{k+1}=x' \mid x_k,u_k].
\]
In words, the distribution of the next state depends only on the current state and current
action, not on how the system reached the current state.

The immediate cost is a function
\[
c:\cX \times \cU \times \cX \to \R,
\]
where
\[
c(x,u,x')
\]
is the cost incurred when the system is in state \(x\), action \(u\) is chosen, and the next state
turns out to be \(x'\). Since \(x'\) is random, it is useful to define the expected one-stage cost
\[
C_x^u
:=
\E[c(x,u,x_{k+1}) \mid x_k=x,\ u_k=u]
=
\sum_{x'\in\cX} P^u_{x,x'} c(x,u,x').
\]

\begin{theorembox}[Finite MDP]
A finite Markov Decision Process is described by
\[
(\cX,\cU,P,c),
\]
where (\cX) is a finite state space, \(\cU\) is a finite action space, \((P^u_{x,x'})\) are the
transition probabilities, and \((c(x,u,x'))\) is the one-step cost.
\end{theorembox}

\subsubsection{Policies}

The control design problem for an MDP consists in selecting actions based on the current state.
This is done through a policy. A \textbf{deterministic policy} is a map
\[
\mu:\cX \to \cU.
\]
If the system is in state \(x\), the controller applies
\[
u = \mu(x).
\]
Thus a deterministic policy assigns exactly one action to each state.\\

A \textbf{stochastic policy}, also called a mixed policy, assigns a probability distribution over actions:
\[
\pi:\cX \times \cU \to [0,1],
\]
where
\[
\pi(x,u) = \Prob[u_k=u \mid x_k=x].
\]
For every state \(x\),
\[
\pi(x,u)\ge 0,
\qquad
\sum_{u\in\cU}\pi(x,u)=1.
\]
Thus, when the system is in state \(x\), the action is sampled according to the probabilities
\(\pi(x,u)\).

Deterministic policies are a special case of stochastic policies. Indeed, a deterministic policy
\(\mu\) corresponds to the stochastic policy
\[
\pi(x,u)
=
\begin{cases}
1, & u=\mu(x),\
0, & u\neq \mu(x).
\end{cases}
\]
In many control applications, one ultimately wants a deterministic controller. Nevertheless, the
stochastic formulation is very convenient mathematically, because it allows us to write policy
optimization problems as optimization problems over probability distributions. This viewpoint
will also be important when discussing learning-based methods.

\subsection{Finite-horizon dynamic programming}
\subsubsection{Finite-horizon sequential decision problems}

We now consider an MDP evolving through time. The index \(k\) denotes the time step. At time
\(k\), the system is in state \(x_k\in\cX\), the controller selects an action \(u_k\in\cU\), and
the next state \(x_{k+1}\) is generated according to the transition probabilities
\[
\Prob[x_{k+1}=x' \mid x_k=x,\ u_k=u] = P^u_{x,x'}.
\]
The one-stage cost incurred at time \(k\) is
\[
c_k = c(x_k,u_k,x_{k+1}).
\]
Thus \(c_k\) is the cost paid when the system moves from \(x_k\) to \(x_{k+1}\) under the
input \(u_k\).

As control designers, we are interested in choosing a policy that gives good performance over a
time horizon. In a finite-horizon problem, a natural performance index is
\[
\sum_{k=0}^{K} c_k.
\]
Since the system is stochastic, this quantity is random. Therefore, for a policy sequence
\[
\pi_0,\pi_1,\dots,\pi_K,
\]
we evaluate performance through the conditional expected cost
\[
V^\pi(x)
=
\E\left[
\sum_{k=0}^{K} c_k
\ \middle|\ x_0=x
\right].
\]
The goal is to find a policy sequence that minimizes this expected cumulative cost.

\subsubsection{Dynamic programming recursion for a fixed policy}

To solve the finite-horizon problem recursively, we introduce the value function from an
intermediate time \(k\):
\[
V_k^\pi(x)
=
\E\left[
\sum_{i=k}^{K} c_i
\ \middle|\ x_k=x
\right].
\]
This quantity represents the expected future cost from time \(k\) to the end of the horizon,
assuming that the system is currently in state \(x\) and that the policy sequence
\[
\pi_k,\pi_{k+1},\dots,\pi_K
\]
is used from time \(k\) onward.

The key point is that the value function can be written recursively. Indeed,
\[
V_k^\pi(x)
=
\E\left[
\sum_{i=k}^{K} c_i
\ \middle|\ x_k=x
\right]
\]
can be split into the immediate cost at time \(k\) and the future cost from time \(k+1\):
\[
V_k^\pi(x)
=
\E\left[
c_k+\sum_{i=k+1}^{K} c_i
\ \middle|\ x_k=x
\right].
\]
Now we condition on the action \(u\) chosen at state \(x\), and on the next state \(x'\). The
policy chooses action \(u\) with probability \(\pi_k(x,u)\), and the system moves to \(x'\) with
probability \(P^u_{x,x'}\). Therefore,
\[
V_k^\pi(x)
=
\sum_{u\in\cU}\pi_k(x,u)
\left[
C_x^u
+
\sum_{x'\in\cX}P^u_{x,x'}
\E\left[
\sum_{i=k+1}^{K}c_i
\ \middle|\ x_{k+1}=x'
\right]
\right],
\]
where
\[
C_x^u
=
\E[c_k\mid x_k=x,\ u_k=u]
\]
is the expected immediate cost when action \(u\) is applied in state \(x\).

By definition of the value function at the next time step, we have
\[
\E\left[
\sum_{i=k+1}^{K}c_i
\ \middle|\ x_{k+1}=x'
\right]
=
V_{k+1}^\pi(x').
\]
Hence the recursion becomes
\[
V_k^\pi(x)
=
\sum_{u\in\cU}\pi_k(x,u)
\left[
C_x^u
+
\sum_{x'\in\cX}P^u_{x,x'}V_{k+1}^\pi(x')
\right].
\]
This is the dynamic programming recursion for the value of a fixed policy.

The recursion should be read backward in time. At the final stage, only the last stage cost
remains. Equivalently, one can set
\[
V_{K+1}^\pi(x)=0
\]
and apply the same recursion for \(k=K,K-1,\dots,0\). Thus the finite-horizon value function is
computed by backward induction.

\subsubsection{Bellman's optimality principle and backward induction}

The optimal value from stage \(k\) is defined as the minimum expected cost that can be achieved
from state \(x\):
\[
V_k^\star(x)
=
\min_{\pi_k,\pi_{k+1},\dots,\pi_K}
\E\left[
\sum_{i=k}^{K}c_i
\ \middle|\ x_k=x
\right].
\]
By Bellman's optimality principle, once the first action has been selected and the next state has
been reached, the remaining decisions must be optimal for the subproblem starting from that
new state. Therefore the optimal value satisfies
\[
V_k^\star(x)
=
\min_{\pi_k(x,\cdot)}
\sum_{u\in\cU}\pi_k(x,u)
\left[
C_x^u
+
\sum_{x'\in\cX}P^u_{x,x'}V_{k+1}^\star(x')
\right].
\]
This equation is the stochastic analogue of the deterministic Bellman recursion. The term
\(C_x^u\) is the immediate expected cost, while
\[
\sum_{x'\in\cX}P^u_{x,x'}V_{k+1}^\star(x')
\]
is the expected optimal future cost after applying action \(u\).

For each fixed state \(x\), the minimization is over the probability distribution
\(\pi_k(x,\cdot)\). Hence it is a linear program over the probability simplex
\[
\Delta(\cU)
=
\left\{
\nu\in\R^M:
\nu_u\ge 0,\ 
\sum_{u\in\cU}\nu_u=1
\right\}.
\]
Indeed, for fixed \(V_{k+1}^\star\), the quantity
\[
C_x^u+\sum_{x'\in\cX}P^u_{x,x'}V_{k+1}^\star(x')
\]
is just a scalar number associated with action \(u\). Therefore the objective is linear in the
probabilities \(\pi_k(x,u)\).

Since a linear function over a simplex attains its minimum at an extreme point, under the usual
finite-state and finite-action assumptions there exists an optimal deterministic action. Thus the
Bellman recursion can be written equivalently as
\[
V_k^\star(x)
=
\min_{u\in\cU}
\left[
C_x^u
+
\sum_{x'\in\cX}P^u_{x,x'}V_{k+1}^\star(x')
\right].
\]
An optimal deterministic policy is then obtained by selecting
\[
\mu_k^\star(x)
\in
\argmin_{u\in\cU}
\left[
C_x^u
+
\sum_{x'\in\cX}P^u_{x,x'}V_{k+1}^\star(x')
\right].
\]

\begin{theorembox}[Finite-horizon Bellman recursion]
For a finite-horizon MDP, set \(V_{K+1}^\star(x)=0\). Then, for
\(k=K,K-1,\dots,0\),
\[
V_k^\star(x)
=
\min_{\pi_k(x,\cdot)}
\sum_{u\in\cU}\pi_k(x,u)
\left[
C_x^u
+
\sum_{x'\in\cX}P^u_{x,x'}V_{k+1}^\star(x')
\right].
\]
Equivalently, since an optimal deterministic action exists,
\[
V_k^\star(x)
=
\min_{u\in\cU}
\left[
C_x^u
+
\sum_{x'\in\cX}P^u_{x,x'}V_{k+1}^\star(x')
\right].
\]
\end{theorembox}

The finite-horizon problem is therefore solved by backward induction:
\begin{itemize}
    \item \(V_{K+1}^\star\) acts as the base case;
    \item \(K+1\) backward steps are performed;
    \item at each step, one policy \(\pi_k\) is found;
    \item for each state \(x\), the backward step is a linear program over the action
    probabilities;
    \item under weak conditions, the optimal policy can be chosen deterministic.
\end{itemize}

This is a very general method: the state space is discrete, the dynamics may be nonlinear, and
the evolution is stochastic. The catch is computational. If the state space is very large, then
the value \(V_k^\star(x)\) must be computed and stored for many states, and this quickly becomes
expensive. This is the usual curse of dimensionality of dynamic programming.

\subsection{Infinite-horizon discounted MDPs}
\subsubsection{Stationary problems and discounted cost}

In many MDP problems, there is no natural final time. For instance, a traffic light controller is
not designed to operate only until a fixed terminal time; it should run continuously. This leads
to infinite-horizon problems.

In the infinite-horizon setting, we typically focus on stationary problems:
\begin{itemize}
    \item the transition probabilities do not depend on time;
    \item the expected cost \(C_x^u\) does not depend on time;
    \item the policy \(\pi\) is time invariant.
\end{itemize}
Thus the same policy is used at every time step:
\[
\pi(x,u)=\Prob[u_k=u\mid x_k=x],
\qquad k=0,1,2,\dots
\]

A standard infinite-horizon performance index is the discounted cost
\[
J
=
\sum_{k=0}^{\infty}\gamma^k c_k,
\]
where
\[
0\le \gamma <1
\]
is the discount factor.

\subsubsection{Discount factor}

The discount factor has several interpretations.

From a mathematical perspective, the condition \(\gamma<1\) ensures that the infinite sum is
finite, provided that the one-stage costs remain bounded. Indeed, the weights
\[
1,\gamma,\gamma^2,\dots
\]
form a convergent geometric sequence.

In some applications, the discount factor represents the lower impact of future costs. For
example, in economics, money spent later may be considered less costly than money spent today,
because money available today could be invested.

Another interpretation is probabilistic. The factor \(\gamma\) can be seen as a proxy for a
Bernoulli termination probability. If at each time step the process continues with probability
\(\gamma\), then future costs are weighted by the probability that the process has not yet
terminated.

The extreme cases are also informative:
\begin{itemize}
    \item if \(\gamma=0\), only the immediate cost is considered;
    \item if \(0<\gamma<1\), the controller balances immediate cost and future cost;
    \item if \(\gamma\) is close to \(1\), future costs become almost as important as current costs.
\end{itemize}

\subsubsection{Closed-loop MDP and stationary distribution}

Once a policy \(\pi\) is fixed, the MDP becomes a Markov chain. Indeed, the action is no longer
a free decision variable: it is sampled according to the policy. Therefore, the transition
probability from \(x\) to \(x'\) under the closed-loop policy is
\[
P^\pi_{x,x'}
=
\sum_{u\in\cU}
\Prob[x'\mid x,u]\Prob[u\mid x]
=
\sum_{u\in\cU}
\pi(x,u)P^u_{x,x'}.
\]
The matrix \(P^\pi\) is the transition matrix of the closed-loop Markov chain.

Let \(d_k(x)\) be the probability that the system is in state \(x\) at time \(k\). If \(d_k\) is
written as a column vector, then the distribution evolves as
\[
d_{k+1}
=
(P^\pi)^\top d_k.
\]
Equivalently, if distributions are written as row vectors, then
\[
d_{k+1}^\top
=
d_k^\top P^\pi.
\]
This evolution is linear in the probability distribution.

A stationary distribution for the closed-loop Markov chain is a probability vector
\(\bar d_\pi\) satisfying
\[
\bar d_\pi^\top
=
\bar d_\pi^\top P^\pi.
\]
Thus \(\bar d_\pi\) is a left eigenvector of \(P^\pi\) associated with the eigenvalue \(1\),
normalized so that its entries sum to one. The component \(\bar d_\pi(x)\) gives the steady-state
probability that the system is in state \(x\).

If the Markov chain is ergodic, meaning that every state can be reached from every other state
through some path, then this stationary distribution is unique and describes the long-run
fraction of time spent in each state.

\subsubsection{Closed-loop queue example}

Consider again the traffic-light queue example, with states
\[
\cX=\{0,1,2,3\}.
\]
Suppose we use the deterministic policy
\[
\mu(q)
=
\begin{cases}
\text{green} & \text{if } q\ge 3,\\
\text{red} & \text{otherwise}.
\end{cases}
\]
Hence the light is red for \(q=0,1,2\), and it turns green when the queue reaches length \(3\).

Using the state ordering \(0,1,2,3\), the closed-loop transition matrix is
\[
P^\pi
=
\begin{bmatrix}
1-p & p   & 0   & 0\\
0   & 1-p & p   & 0\\
0   & 0   & 1-p & p\\
1-p & p   & 0   & 0
\end{bmatrix}.
\]
The first row means that, from queue length \(0\), the queue stays at \(0\) with probability
\(1-p\) and moves to \(1\) with probability \(p\). The same logic applies to queue lengths
\(1\) and \(2\). The last row is different because, when \(q=3\), the light is green and three
cars are served; after that, either no new car arrives and the next queue length is \(0\), or one
new car arrives and the next queue length is \(1\).

The stationary distribution is obtained by solving
\[
\bar d_\pi^\top
=
\bar d_\pi^\top P^\pi
\]
together with the normalization condition
\[
\sum_{x\in\cX}\bar d_\pi(x)=1.
\]
For example, if \(p=0.2\), one obtains approximately
\[
\bar d_\pi
=
\begin{bmatrix}
0.267\\
0.333\\
0.333\\
0.067
\end{bmatrix}.
\]
Therefore, under this policy, the queue is empty about \(26.7\%\) of the time, has one car about
\(33.3\%\) of the time, has two cars about \(33.3\%\) of the time, and has three cars about
\(6.7\%\) of the time.

\subsection{Policy evaluation and Bellman optimality}
\subsubsection{Infinite-horizon value of a policy}

For a fixed stationary policy \(\pi\), the infinite-horizon discounted value is defined as
\[
V^\pi(x)
=
\E\left[
\sum_{k=0}^{\infty}\gamma^k c_k
\ \middle|\ x_0=x
\right].
\]
This is the conditional expected future cost when the process starts from state \(x\) and policy
\(\pi\) is used forever.

The value of a policy satisfies the discounted Bellman equation
\[
V^\pi(x)
=
\sum_{u\in\cU}\pi(x,u)
\left[
C_x^u
+
\gamma
\sum_{x'\in\cX}P^u_{x,x'}V^\pi(x')
\right].
\]
To see why, split the infinite-horizon sum into the immediate cost and the discounted future
cost:
\[
V^\pi(x)
=
\E\left[
c_0+\gamma\sum_{k=1}^{\infty}\gamma^{k-1}c_k
\ \middle|\ x_0=x
\right].
\]
After the first transition, the expected remaining discounted cost is exactly \(V^\pi(x_1)\).
Conditioning on the action \(u\) and on the next state \(x'\) gives the Bellman equation above.

This equation is a consistency equation. It says that the value of a state must be equal to the
expected immediate cost plus the discounted value of the next state. The system dynamics are
encoded in \(P^u_{x,x'}\), while the cost function is encoded in \(C_x^u\).

This interpretation is particularly important in learning. The value \(V^\pi(x)\) is a predictor
of the quality of a candidate policy. The Bellman equation decomposes this predictor into two
parts:
\begin{itemize}
    \item the immediate observable cost \(C_x^u\);
    \item the estimate of future cost \(V^\pi(x')\).
\end{itemize}

\subsubsection{Computational complexity of policy evaluation}

For a fixed policy \(\pi\), define the expected cost vector
\[
C^\pi(x)
=
\sum_{u\in\cU}\pi(x,u)C_x^u
\]
and the closed-loop transition matrix
\[
P^\pi_{x,x'}
=
\sum_{u\in\cU}\pi(x,u)P^u_{x,x'}.
\]
Then the discounted Bellman equation can be written in vector form as
\[
V^\pi
=
C^\pi+\gamma P^\pi V^\pi.
\]
Rearranging gives
\[
(I-\gamma P^\pi)V^\pi
=
C^\pi,
\]
and therefore
\[
V^\pi
=
(I-\gamma P^\pi)^{-1}C^\pi.
\]

\begin{theorembox}[Policy evaluation]
For a fixed stationary policy \(\pi\), the value function is the unique solution of
\[
V^\pi
=
C^\pi+\gamma P^\pi V^\pi.
\]
Equivalently,
\[
V^\pi
=
(I-\gamma P^\pi)^{-1}C^\pi.
\]
\end{theorembox}

The matrix \(I-\gamma P^\pi\) is invertible for \(\gamma\in[0,1)\). Intuitively, the eigenvalues
of \(P^\pi\) lie in the unit disk for a stochastic matrix, and multiplying by \(\gamma<1\) moves
them strictly inside the unit disk. Thus \(1\) cannot be an eigenvalue of \(\gamma P^\pi\), and
\(I-\gamma P^\pi\) is nonsingular.

If the MDP has \(N\) states, the Bellman equation for a fixed policy is a system of \(N\)
linear equations. Solving it directly can be computationally expensive for large systems. In a
dense representation, the cost is on the order of
\[
O(N^3+N^2M),
\]
where the \(O(N^3)\) term comes from solving the linear system and the \(O(N^2M)\) term comes
from constructing the closed-loop quantities from all state-action transition matrices.

As an alternative, one can iteratively update the Bellman equation. This is possible because the
discounted Bellman operator for a fixed policy is contractive when \(\gamma<1\).

\subsubsection{Queue example for policy evaluation}

For the queue example with the policy described above, the closed-loop matrix is
\[
P^\pi
=
\begin{bmatrix}
1-p & p   & 0   & 0\\
0   & 1-p & p   & 0\\
0   & 0   & 1-p & p\\
1-p & p   & 0   & 0
\end{bmatrix}.
\]
If the cost is the number of cars in the queue, then
\[
C^\pi
=
\begin{bmatrix}
0\\
1\\
2\\
3
\end{bmatrix}.
\]
The value of the policy is therefore
\[
V^\pi
=
(I-\gamma P^\pi)^{-1}C^\pi.
\]
This gives, for each initial queue length, the expected discounted total number of cars in the
queue when the policy is used forever.

\subsubsection{Bellman principle for infinite-horizon MDPs}

The optimal value is defined as the best value achievable among all policies:
\[
V^\star(x)
=
\min_{\textcolor{red}{\pi}} V^{\textcolor{red}{\pi}} (x).
\]
Using the Bellman equation for a policy, we can write
\[
V^\star(x)
=
\min_{\textcolor{red}{\pi}}
\sum_{u\in\cU}\textcolor{red}{\pi(x,u)}
\left[
C_x^u
+
\gamma\sum_{x'\in\cX}P^u_{x,x'}V^{\textcolor{red}{\pi}}(x')
\right].
\]
Bellman's optimality principle turns this into a recursive equation for \(V^\star\):
\[
\textcolor{blue}{V^\star(x)}
=
\min_{\textcolor{red}{\pi}}
\sum_{u\in\cU}\textcolor{red}{\pi(x,u)}
\left[
C_x^u
+
\gamma\sum_{x'\in\cX}P^u_{x,x'}\textcolor{blue}{V^\star(x')}
\right].
\]
As before, the minimization is linear in the action probabilities \(\pi(x,u)\). Therefore, if the
minimum is achieved at a deterministic policy, the equation is equivalent to
\[
\textcolor{blue}{V^\star(x)}
=
\min_{\textcolor{red}{u}\in\cU}
\left[
C_x^{\textcolor{red}{u}}
+
\gamma\sum_{x'\in\cX}P^{\textcolor{red}{u}}_{x,x'}\textcolor{blue}{V^\star(x')}
\right].
\]
This is a system of \(N\) nonlinear equations. The nonlinearity does not come from nonlinear
dynamics, but from the minimum operator over the finite action set.

\begin{theorembox}[Infinite-horizon Bellman optimality equation]
The optimal value function of a discounted finite MDP satisfies
\[
V^\star(x)
=
\min_{u\in\cU}
\left[
C_x^u
+
\gamma\sum_{x'\in\cX}P^u_{x,x'}V^\star(x')
\right],
\qquad x\in\cX.
\]
An optimal deterministic policy is obtained by choosing
\[
\mu^\star(x)
\in
\argmin_{u\in\cU}
\left[
C_x^u
+
\gamma\sum_{x'\in\cX}P^u_{x,x'}V^\star(x')
\right].
\]
\end{theorembox}

\subsection{Iterative computation of the optimal policy}
\subsubsection{Greedy policy improvement}
The Bellman optimality equation characterizes the optimal value function, but it does not
immediately give a practical way to compute the optimal policy. We now discuss two classical
iterative methods for doing this: policy iteration and value iteration.

The first ingredient is \textbf{policy evaluation}. Given a policy \(\pi\), we can compute the value
function \(V^\pi\) associated with that policy by solving the discounted Bellman equation
\[
V^\pi
=
C^\pi+\gamma P^\pi V^\pi.
\]
Equivalently,
\[
V^\pi
=
(I-\gamma P^\pi)^{-1}C^\pi.
\]
This step tells us how good the current policy is.

At the same time, once we know the value \(V^\pi\), we can try to improve the policy. The idea
is to ask the following question: if the system is currently in state \(x\), can we reduce the cost
by deviating from the current policy for one step, and then following the old policy again from
the next state onward?

For a candidate stochastic action distribution \(\nu(x,\cdot)\in\Delta(\cU)\), the one-step
lookahead cost is
\[
\sum_{u\in\cU}\nu(x,u)
\left[
C_x^u
+
\gamma\sum_{x'\in\cX}P^u_{x,x'}V^\pi(x')
\right].
\]
Therefore, the improved policy is obtained by solving
\[
\pi'(x,\cdot)
\in
\argmin_{\nu(x,\cdot)\in\Delta(\cU)}
\sum_{u\in\cU}\nu(x,u)
\left[
C_x^u
+
\gamma\sum_{x'\in\cX}P^u_{x,x'}V^\pi(x')
\right].
\]
Since this minimization is linear over the probability simplex, the optimizer can be chosen as
a deterministic policy. Hence we can equivalently write
\[
\mu'(x)
\in
\argmin_{u\in\cU}
\left[
C_x^u
+
\gamma\sum_{x'\in\cX}P^u_{x,x'}V^\pi(x')
\right].
\]
The interpretation is simple: at each state \(x\), we choose the action that gives the smallest
sum of immediate expected cost and discounted future value, assuming that after this first
action we continue with the old policy \(\pi\).

\subsubsection{Policy improvement theorem}

The greedy improvement step always improves the value, unless the current policy is already
optimal. More precisely, let \(\pi'\) be the greedy policy obtained from \(V^\pi\). Then
\[
V^{\pi'}(x)
\le
V^\pi(x)
\qquad
\forall x\in\cX.
\]
Moreover, unless \(V^\pi\) is already the optimal value, there exists at least one state \(x'\)
such that
\[
V^{\pi'}(x')
<
V^\pi(x').
\]

\begin{theorembox}[Policy improvement theorem]
Let \(\pi\) be a policy and let \(\pi'\) be obtained by the greedy one-step improvement with
respect to \(V^\pi\). Then
\[
V^{\pi'}(x)\le V^\pi(x)
\qquad
\forall x\in\cX.
\]
If the greedy improvement is strict for some state, then the new policy is strictly better for
at least one state. If no strict improvement is possible, then \(\pi\) is optimal.
\end{theorembox}

\paragraph{Proof.}
We prove the inequality
\[
V^{\pi'}(x)\le V^\pi(x).
\]
By definition of the greedy policy \(\pi'\), for every state \(x\) we have
\[
V^\pi(x)
\ge
\sum_{u\in\cU}\pi'(x,u)
\left[
C_x^u
+
\gamma\sum_{x'\in\cX}P^u_{x,x'}V^\pi(x')
\right].
\]
The right-hand side is the expected value of
\[
c_0+\gamma V^\pi(x_1)
\]
when the first action is chosen according to \(\pi'\). Hence
\[
V^\pi(x)
\ge
\E_{\pi'}\left[
c_0+\gamma V^\pi(x_1)
\ \middle|\ x_0=x
\right].
\]
Now we apply the same inequality again from the next state \(x_1\). This gives
\[
V^\pi(x)
\ge
\E_{\pi'}\left[
c_0+\gamma\left(c_1+\gamma V^\pi(x_2)\right)
\ \middle|\ x_0=x
\right],
\]
or equivalently
\[
V^\pi(x)
\ge
\E_{\pi'}\left[
c_0+\gamma c_1+\gamma^2V^\pi(x_2)
\ \middle|\ x_0=x
\right].
\]
Repeating the same argument recursively, we obtain
\[
V^\pi(x)
\ge
\E_{\pi'}\left[
c_0+\gamma c_1+\gamma^2c_2+\gamma^3c_3+\cdots
\ \middle|\ x_0=x
\right].
\]
The right-hand side is precisely the value of the improved policy:
\[
\E_{\pi'}\left[
c_0+\gamma c_1+\gamma^2c_2+\gamma^3c_3+\cdots
\ \middle|\ x_0=x
\right]
=
V^{\pi'}(x).
\]
Therefore
\[
V^\pi(x)\ge V^{\pi'}(x)
\qquad
\forall x\in\cX.
\]
The strict inequality follows if the greedy step gives a strict improvement for at least one
state, while being non-worse for all the others.

As a consequence, after a finite number of strict policy improvements, we reach an optimal
policy. The reason is that there are only finitely many deterministic policies in a finite MDP,
and policy improvement never returns to a strictly worse policy.

\subsubsection{Policy iteration algorithm}

Policy iteration alternates between the two operations described above:
\begin{itemize}
    \item evaluate the current policy;
    \item greedily improve the current policy using the value function just computed.
\end{itemize}

\begin{theorembox}[Policy iteration algorithm]
Initialize a policy guess
\[
\pi \leftarrow \pi_0.
\]
For each iteration:
\begin{enumerate}
    \item \textbf{Policy evaluation.} Compute the value \(V\) associated with the current policy
    \(\pi\):
    \[
    V
    \leftarrow
    (I-\gamma P^\pi)^{-1}C^\pi.
    \]

    \item \textbf{Greedy policy update.} For every state \(x\), update the policy by solving
    \[
    \pi_{\mathrm{new}}(x,\cdot)
    \in
    \argmin_{\nu(x,\cdot)\in\Delta(\cU)}
    \sum_{u\in\cU}\nu(x,u)
    \left[
    C_x^u
    +
    \gamma\sum_{x'\in\cX}P^u_{x,x'}V(x')
    \right].
    \]
    Equivalently, since the minimizer can be chosen deterministic,
    \[
    \mu_{\mathrm{new}}(x)
    \in
    \argmin_{u\in\cU}
    \left[
    C_x^u
    +
    \gamma\sum_{x'\in\cX}P^u_{x,x'}V(x')
    \right].
    \]

    \item Set
    \[
    \pi \leftarrow \pi_{\mathrm{new}}.
    \]
\end{enumerate}
Stop when the policy no longer changes.
\end{theorembox}

The computational cost is split between the two steps. Policy evaluation is the expensive part,
because it requires solving a system of \(N\) linear equations:
\[
V=(I-\gamma P^\pi)^{-1}C^\pi.
\]
For a dense MDP, this has complexity approximately
\[
O(N^3+N^2M).
\]
The term \(O(N^3)\) comes from solving the linear system, while the term \(O(N^2M)\) comes
from constructing the closed-loop transition matrix and expected cost.

Policy improvement is usually simpler. For each state \(x\), one compares the \(M\) possible
actions and selects the one with smallest one-step lookahead cost. Thus the improvement step is
essentially a minimization over \(M\) alternatives, repeated \(N\) times.

The policy improvement theorem guarantees that every nontrivial iteration improves the policy.
Since the number of deterministic policies is finite, policy iteration terminates after a finite
number of improvements.

\subsubsection{Value iteration}

Policy iteration evaluates each intermediate policy exactly. An alternative is to look directly
for a fixed point of the Bellman optimality equation.

Recall that the optimal value satisfies
\[
V^\star(x)
=
\min_{\pi(x,\cdot)}
\sum_{u\in\cU}\pi(x,u)
\left[
C_x^u
+
\gamma\sum_{x'\in\cX}P^u_{x,x'}V^\star(x')
\right].
\]
Equivalently, when the optimum is attained by a deterministic policy,
\[
V^\star(x)
=
\min_{u\in\cU}
\left[
C_x^u
+
\gamma\sum_{x'\in\cX}P^u_{x,x'}V^\star(x')
\right].
\]

A natural way to compute this fixed point is to start from an initial guess \(V^{(0)}\) and
iterate the Bellman optimality equation:
\[
V^{(t+1)}(x)
=
\min_{\pi(x,\cdot)}
\sum_{u\in\cU}\pi(x,u)
\left[
C_x^u
+
\gamma\sum_{x'\in\cX}P^u_{x,x'}V^{(t)}(x')
\right],
\]
or, equivalently,
\[
V^{(t+1)}(x)
=
\min_{u\in\cU}
\left[
C_x^u
+
\gamma\sum_{x'\in\cX}P^u_{x,x'}V^{(t)}(x')
\right].
\]
Here \(t\) is not the time index of the MDP. It is the iteration index of the algorithm.

To state the convergence property compactly, define the Bellman optimality operator \(T\) as
\[
(TV)(x)
=
\min_{u\in\cU}
\left[
C_x^u
+
\gamma\sum_{x'\in\cX}P^u_{x,x'}V(x')
\right].
\]
Then value iteration is simply
\[
V^{(t+1)}=TV^{(t)}.
\]

The key convergence result is that the Bellman optimality operator is a contraction in the
infinity norm. If \(V\) and \(W\) are two value functions, then
\[
\norm{TV-TW}_{\infty}
\le
\gamma\norm{V-W}_{\infty},
\]
where
\[
\norm{V}_{\infty}
:=
\max_{x\in\cX}|V(x)|.
\]
Since \(0\le\gamma<1\), the operator \(T\) is contractive. Therefore, value iteration converges
to the unique fixed point \(V^\star\). In particular,
\[
\norm{V^{(t)}-V^\star}_{\infty}
\le
\gamma^t\norm{V^{(0)}-V^\star}_{\infty}.
\]
Thus the convergence rate is governed by \(\gamma\). When \(\gamma\) is close to one, future
costs are important, but the convergence of value iteration becomes slower.

\subsubsection{Value iteration algorithm}

\begin{theorembox}[Value iteration algorithm]
Initialize a value guess
\[
V\leftarrow V_0.
\]
For each iteration, apply the Bellman iteration
\[
V(x)
\leftarrow
\min_{\pi(x,\cdot)}
\sum_{u\in\cU}\pi(x,u)
\left[
C_x^u
+
\gamma\sum_{x'\in\cX}P^u_{x,x'}V(x')
\right].
\]
Equivalently,
\[
V(x)
\leftarrow
\min_{u\in\cU}
\left[
C_x^u
+
\gamma\sum_{x'\in\cX}P^u_{x,x'}V(x')
\right],
\qquad x\in\cX.
\]

At convergence, extract the optimal greedy policy:
\[
\pi^\star(x,\cdot)
\in
\argmin_{\nu(x,\cdot)\in\Delta(\cU)}
\sum_{u\in\cU}\nu(x,u)
\left[
C_x^u
+
\gamma\sum_{x'\in\cX}P^u_{x,x'}V^\star(x')
\right].
\]
Equivalently, one can extract a deterministic optimal policy by choosing
\[
\mu^\star(x)
\in
\argmin_{u\in\cU}
\left[
C_x^u
+
\gamma\sum_{x'\in\cX}P^u_{x,x'}V^\star(x')
\right].
\]
\end{theorembox}

The computational cost of value iteration is different from policy iteration. Each Bellman
iteration is usually cheaper than solving the full policy evaluation linear system. For a dense
finite MDP, one Bellman iteration has complexity approximately
\[
O(N^2M).
\]
However, the convergence is asymptotic rather than finite, and the rate is controlled by
\(\gamma\). Thus policy iteration usually performs fewer but more expensive iterations, while
value iteration performs cheaper iterations but may need many of them.